%% Chapter 8 — Section 8.3: Discussion and Bibliography
%% Source: PDF page 108 (book page 95)

\section{Discussion and Bibliography}
\label{sec:8.3}

HMC was proposed in the physics literature under the name Hybrid Monte Carlo~\cite{60}. The
method was introduced to the statistics community through the work~\cite{171}. An accessible
introduction to HMC is~\cite{170}; see also~\cite{23}. Theorem~8.4, which shows that
Hamiltonian flows are symplectic, was proved in~\cite{187}. Further background on Hamiltonian
dynamics and their numerical solution can be found in~\cite{210,101}. Numerical integrators
tailored to HMC are reviewed and compared in~\cite{31,42}.

In this chapter, we have seen that the Metropolis Hastings framework can be generalized by
means of an involution map, which we take as a momentum flip for HMC methods. The use of
involutions within the Metropolis Hastings algorithm was proposed in~\cite{11,92}, and has been
leveraged to establish mixing rates for Hamiltonian Monte Carlo in~\cite{92}. The theory
in~\cite{92}, which relies on the weak Harris approach developed in~\cite{102}, also covers
implementations of HMC in infinite-dimensional Hilbert spaces~\cite{22}. The paper~\cite{30}
relies instead on a novel coupling technique to analyze the convergence of HMC algorithms.
Further convergence results can be found in~\cite{148,65} and optimal scaling of HMC in high
dimensions was studied in~\cite{21}. Tuning HMC algorithms can be challenging. The No-U-Turn
Sampler (NUTS) algorithm~\cite{113} introduces a computational framework to automatically set
the duration parameter $\lambda$ and the step-size $\epsilon$.
